\documentclass[11pt]{article}

\usepackage[letterpaper,top=1in,bottom=1in,left=1in,right=1in]{geometry}
\usepackage{amsmath,amsthm,amssymb}
\usepackage{hyperref}

\newtheorem{theorem}{Theorem}[section]
\newtheorem{lemma}[theorem]{Lemma}
\newtheorem{proposition}[theorem]{Proposition}
\newtheorem{corollary}[theorem]{Corollary}
\theoremstyle{definition}
\newtheorem{definition}[theorem]{Definition}
\newtheorem{example}[theorem]{Example}
\theoremstyle{remark}
\newtheorem{remark}[theorem]{Remark}
\newtheorem{question}[theorem]{Question}

\newcommand{\R}{\mathbb{R}}
\newcommand{\Q}{\mathbb{Q}}
\newcommand{\N}{\mathbb{N}}
\newcommand{\bM}{\bar{\mathcal{M}}}
\newcommand{\ex}{\operatorname{ex}}
\newcommand{\eps}{\varepsilon}
\newcommand{\E}{\mathrm{E}}
\newcommand{\GI}{\cong_{\mathrm{gi}}}
\newcommand{\isoi}{\cong_i}
\newcommand{\isox}{\cong_i^{\mathrm{ext}}}
\newcommand{\leB}{\leq_B}

\title{\textbf{Every Separable Metric Space Extends to an Extreme One}}

\author{David Victor Feldman\\
\small Department of Mathematics and Statistics\\
\small University of New Hampshire
\and
Eric R.\ Kehoe\\
\small Zoetis, Inc.\\
\small Fort Collins, Colorado}

\date{}

\begin{document}
\maketitle

\begin{abstract}
The bounded-by-$1$ pseudometrics on a set $X$ form a compact convex set $\bM(X)$. For
finite $X$ this is a rational polytope, so its extreme points have rational distances and
satisfy the rigidity constraints that come with lying on many facets. We show that no such
restriction survives passage to infinite $X$: every separable bounded-by-$1$ pseudometric
space $(X,d)$ is the restriction of an extreme point of $\bM(\tilde X)$ for some
$\tilde X \supseteq X$ with $\tilde X \setminus X$ countable. The local structure of an
extreme point may therefore be arbitrary, irrational distances included. The proof glues
finite extreme metrics of prescribed rational distances, supplied by the bowtie
construction of the companion paper, onto a countable dense subset, after subdividing each
irrational distance into a summable sequence of rational ones.

The extension admits a canonical form, functorial in isometric embeddings of countable
spaces, and the original space is recoverable from it as the derived set of its completion.
It follows that graph isomorphism Borel-reduces to isometry of extreme Polish metric
spaces: extreme spaces admit no classification by real-number invariants, nor by countable
structures. A structural obstruction confines this method below the universal orbit
equivalence relation, and the resulting gap is posed as a question.
\end{abstract}

%%%%%%%%%%%%%%%%%%%%%%%%%%%%%%%%%%%%%%%%%%%%%%%%%%%%%%%%%%
\section{Introduction}
\label{sec:intro}
%%%%%%%%%%%%%%%%%%%%%%%%%%%%%%%%%%%%%%%%%%%%%%%%%%%%%%%%%%

Let $X$ be a set and let $\bM(X)$ denote the set of bounded-by-$1$ pseudometrics on $X$:
symmetric functions $d\colon X\times X\to[0,1]$ vanishing on the diagonal and satisfying
the triangle inequality. Under pointwise operations $\bM(X)$ is convex, and in the product
topology it is compact, being a closed subset of $[0,1]^{X\times X}$. Write $\ex(\bM(X))$
for its extreme points.

For finite $X$ with $|X|=n$, $\bM(X)$ is a polytope in $\R^{\binom{n}{2}}$ cut out by the
triangle inequalities and the bounds $0\le d\le 1$, all of which have integer coefficients.
Its extreme points are therefore rational, and each is pinned by $\binom{n}{2}$
independent tight constraints---a strong rigidity condition. Which finite metrics are
extreme is a delicate question with a substantial literature
\cite{Avis80,Bandelt92,Deza97,Koolen00}, and even the classification of the extreme rays of the
associated cone is complete only for at most six points.

The situation for infinite $X$ is the opposite, and that is the first content of this
paper.

\begin{theorem}\label{thm:main}
Let $(X,d)$ be a separable bounded-by-$1$ pseudometric space. Then there is a set
$\tilde X\supseteq X$ with $\tilde X\setminus X$ countable, and an extreme point
$\tilde d\in\ex(\bM(\tilde X))$, with $\tilde d|_{X\times X}=d$. The space
$(\tilde X,\tilde d)$ is again separable.
\end{theorem}

So extremality imposes no constraint whatever on the isometry type of a separable
bounded-by-$1$ pseudometric space, up to the addition of countably many points: the class
of extreme separable bounded-by-$1$ pseudometric spaces is universal for the class of all
of them. In particular:

\begin{corollary}\label{cor:irrational}
For countable $X$, an extreme point of $\bM(X)$ may take irrational values, and indeed may
realize any prescribed countable subset of $[0,1]$ among its distances.
\end{corollary}

The mechanism is a trade between two kinds of rigidity. A perturbation of a pseudometric
must vanish wherever a distance is $0$ or $1$, and must respect every degenerate triangle;
in a finite space these constraints compete with the $\binom{n}{2}$ degrees of freedom and
are hard to satisfy in bulk. Given countably many new points one may instead \emph{install}
the constraints: every distance one wishes to freeze is embedded in a finite extreme
metric that realizes it, and the finitely many new distances that the embedding creates are
either equal to $1$, hence frozen outright, or forced by a degenerate triangle through the
gluing locus, hence not free either. Irrational distances, which no finite extreme metric
can realize, are first broken into a summable sequence of rational ones.

The second content concerns classification, and begins from an observation of Dmitri
Nikshych: Theorem~\ref{thm:main} reads as a \emph{nonclassification} result, the world of
extreme separable metric spaces being at least as rich as the world of all of them. We make
this precise. The extension admits a canonical form $\E$, functorial in isometric
embeddings of countable metric spaces (Theorem~\ref{thm:functor}), which in particular
resolves the dependence of the construction on its many choices; and $\E$ is reversible:
the input is recoverable from the output as the derived set of its completion
(Theorem~\ref{thm:recovery}). Composing the two yields a Borel reduction of graph
isomorphism to isometry of extreme Polish metric spaces (Theorem~\ref{thm:gi}), so extreme
Polish metric spaces cannot be classified by real-number invariants, nor by countable
structures (Corollary~\ref{cor:nonclass}). Against the Gao--Kechris theorem
\cite{GaoKechris03}, which places isometry of all Polish metric spaces at the maximal
complexity among orbit equivalence relations of Polish group actions, a gap remains; we
prove that no construction canonical merely at the level of a marked countable dense set
can close it (Proposition~\ref{prop:obstruction}), and pose the sharpened question
(Question~\ref{q:universal}).

Section~\ref{sec:prelim} records what a perturbation of a pseudometric can do; the bound
$|\eps|\le\min\{d,1-d\}$ of Lemma~\ref{lem:bound} is used throughout and is what makes
extremality a property detectable on a dense subset (Lemma~\ref{lem:dense}).
Section~\ref{sec:glue} proves the gluing lemma, Section~\ref{sec:subdivide} the subdivision
lemma, and Section~\ref{sec:proof} assembles them. Section~\ref{sec:canonical} constructs
the canonical extension and Section~\ref{sec:nonclass} the recovery and the reduction.
Section~\ref{sec:questions} collects the remaining questions.

We use one input from the finite theory, proved in Section~\ref{sec:proof} as
Proposition~\ref{prop:realize}: every rational number in $(0,1]$ occurs as a distance in
some extreme point of $\bM(F)$ with $F$ finite. It rests on the generalized bowtie metrics
of the companion paper \cite{FK-cone}, and is the only place the finite classification
enters.

The results of Sections~\ref{sec:prelim}--\ref{sec:proof} have been formalized in Lean~4
against Mathlib; formalization of Sections~\ref{sec:canonical}--\ref{sec:nonclass} is in
progress. Section~\ref{sec:verify} records the correspondence, the scope, and the deposit.

%%%%%%%%%%%%%%%%%%%%%%%%%%%%%%%%%%%%%%%%%%%%%%%%%%%%%%%%%%
\section{Perturbations}
\label{sec:prelim}
%%%%%%%%%%%%%%%%%%%%%%%%%%%%%%%%%%%%%%%%%%%%%%%%%%%%%%%%%%

Throughout, a \emph{metric} means a bounded-by-$1$ pseudometric, and $(X,d)$ ranges over
such spaces of arbitrary cardinality. We write $d(x,y)$ and speak of the \emph{distance}
of the pair $\{x,y\}$.

\begin{definition}\label{def:perturbation}
A \emph{perturbation} of $d\in\bM(X)$ is a symmetric function $\eps\colon X\times X\to\R$
vanishing on the diagonal such that both $d+\eps$ and $d-\eps$ lie in $\bM(X)$.
\end{definition}

Perturbations detect extremality, by the same argument as in a finite-dimensional polytope.

\begin{proposition}\label{prop:extreme-iff}
$d\in\ex(\bM(X))$ if and only if the only perturbation of $d$ is $\eps=0$.
\end{proposition}

\begin{proof}
If $\eps\neq0$ is a perturbation then $d=\tfrac12(d+\eps)+\tfrac12(d-\eps)$ exhibits $d$ as
the midpoint of a segment with distinct endpoints in $\bM(X)$, so $d$ is not extreme.
Conversely if $d=\tfrac12(a+b)$ with $a\neq b$ in $\bM(X)$, then $\eps=\tfrac12(a-b)$ is a
nonzero perturbation.
\end{proof}

The next two lemmas are elementary and do all of the analytic work in this paper.

\begin{lemma}\label{lem:bound}
Let $\eps$ be a perturbation of $d\in\bM(X)$. Then
\[
|\eps(x,y)|\;\le\;\min\bigl\{\,d(x,y),\;1-d(x,y)\,\bigr\}
\qquad\text{for all }x,y\in X .
\]
In particular $\eps(x,y)=0$ whenever $d(x,y)\in\{0,1\}$.
\end{lemma}

\begin{proof}
Both $d+\eps$ and $d-\eps$ take values in $[0,1]$. Nonnegativity of $d-\eps$ and of
$d+\eps$ gives $|\eps|\le d$; the bounds $d\pm\eps\le1$ give $|\eps|\le 1-d$.
\end{proof}

\begin{lemma}\label{lem:tight}
Let $\eps$ be a perturbation of $d\in\bM(X)$ and suppose $d(x,z)=d(x,y)+d(y,z)$. Then
\[
\eps(x,z)=\eps(x,y)+\eps(y,z).
\]
\end{lemma}

\begin{proof}
Applying the triangle inequality to $d+\eps$ at $[x,y,z]$ and subtracting the hypothesis
gives $\eps(x,z)\le\eps(x,y)+\eps(y,z)$; the same computation for $d-\eps$ gives the
reverse inequality.
\end{proof}

Lemma~\ref{lem:bound} has a consequence with no finite-dimensional analogue: a perturbation
is continuous with respect to $d$, whether or not $d+\eps$ has any a priori relation to the
topology of $d$. Extremality is therefore visible on any dense subset.

\begin{lemma}\label{lem:dense}
Let $d\in\bM(X)$, let $A\subseteq X$ be dense in the pseudometric topology of $d$, and let
$\eps$ be a perturbation of $d$ with $\eps|_{A\times A}=0$. Then $\eps=0$.
\end{lemma}

\begin{proof}
Write $\rho=d+\eps\in\bM(X)$. Fix $x,y\in X$ and choose sequences $a_k\to x$ and $b_k\to y$
in $A$. By Lemma~\ref{lem:bound}, $\rho(a_k,x)\le d(a_k,x)+|\eps(a_k,x)|\le 2d(a_k,x)\to0$,
and likewise $\rho(b_k,y)\to 0$. The triangle inequality for $\rho$ gives
\[
|\rho(a_k,b_k)-\rho(x,y)|\;\le\;\rho(a_k,x)+\rho(b_k,y)\;\longrightarrow\;0 ,
\]
so $\rho(a_k,b_k)\to\rho(x,y)$. But $\eps(a_k,b_k)=0$, so $\rho(a_k,b_k)=d(a_k,b_k)$, and
$d(a_k,b_k)\to d(x,y)$ by the same estimate applied to $d$. Hence $\rho(x,y)=d(x,y)$, that
is $\eps(x,y)=0$.
\end{proof}

\begin{corollary}\label{cor:restrict}
If $Y\subseteq X$ and $\eps$ is a perturbation of $d\in\bM(X)$, then $\eps|_{Y\times Y}$ is
a perturbation of $d|_{Y\times Y}$. Consequently, if $d|_{Y\times Y}$ is an extreme point of
$\bM(Y)$ then every perturbation of $d$ vanishes on $Y\times Y$.
\end{corollary}

\begin{proof}
The restriction of a bounded-by-$1$ pseudometric to a subset is one, and restriction is
compatible with the operations $d\pm\eps$.
\end{proof}

%%%%%%%%%%%%%%%%%%%%%%%%%%%%%%%%%%%%%%%%%%%%%%%%%%%%%%%%%%
\section{Gluing along a finite set}
\label{sec:glue}
%%%%%%%%%%%%%%%%%%%%%%%%%%%%%%%%%%%%%%%%%%%%%%%%%%%%%%%%%%

Two metric spaces meeting in a finite set on which they agree may be amalgamated, and the
amalgam can be chosen so that every new distance is either maximal or forced.

\begin{lemma}\label{lem:glue}
Let $(X,d)$ and $(Y,\rho)$ be bounded-by-$1$ pseudometric spaces with $F:=X\cap Y$ finite
and non-empty, and with $d|_{F\times F}=\rho|_{F\times F}$. Define
$\omega\colon (X\cup Y)^2\to[0,1]$ by $\omega|_{X\times X}=d$,
$\omega|_{Y\times Y}=\rho$, and
\begin{equation}\label{eq:glue}
\omega(x,y)\;:=\;\min\Bigl\{\;\min_{z\in F}\bigl(d(x,z)+\rho(z,y)\bigr),\;1\;\Bigr\}
\qquad (x\in X\setminus Y,\; y\in Y\setminus X).
\end{equation}
Then $\omega\in\bM(X\cup Y)$, and for every $x\in X\setminus Y$ and $y\in Y\setminus X$
either $\omega(x,y)=1$, or there is $z\in F$ with
\[
\omega(x,y)=\omega(x,z)+\omega(z,y).
\]
\end{lemma}

\begin{proof}
Symmetry, vanishing on the diagonal and the range $[0,1]$ are clear from
\eqref{eq:glue}; the final dichotomy holds because $F$ is finite, so the inner minimum is
attained, and $\omega(x,z)=d(x,z)$, $\omega(z,y)=\rho(z,y)$ for $z\in F$.

For the triangle inequality write $P=X\setminus Y$, $Q=Y\setminus X$. A triple with all
three points in $X$, or all three in $Y$, is governed by $d$ or by $\rho$. Since $F\subseteq
X$ and $F\subseteq Y$, the remaining triples have at least one point in $P$ and at least one
in $Q$, leaving three cases up to the symmetry exchanging $(X,d,P)$ with $(Y,\rho,Q)$.

\smallskip
\noindent\emph{Case 1: $x\in P$, $y\in Q$, $z\in F$.}
First, $\omega(x,y)\le d(x,z)+\rho(z,y)=\omega(x,z)+\omega(z,y)$ directly from
\eqref{eq:glue}. Next, $\omega(x,z)\le\omega(x,y)+\omega(y,z)$: if $\omega(x,y)=1$ this
reads $d(x,z)\le1+\rho(y,z)$, which holds as $d\le1$; otherwise
$\omega(x,y)=d(x,w)+\rho(w,y)$ for some $w\in F$, and
\[
d(x,z)\;\le\;d(x,w)+d(w,z)\;=\;d(x,w)+\rho(w,z)\;\le\;d(x,w)+\rho(w,y)+\rho(y,z),
\]
using $d|_F=\rho|_F$. The inequality $\omega(z,y)\le\omega(x,y)+\omega(x,z)$ follows by
exchanging the roles of $x$ and $y$.

\smallskip
\noindent\emph{Case 2: $x,x'\in P$, $y\in Q$.}
For $\omega(x,y)\le\omega(x,x')+\omega(x',y)$: if $\omega(x',y)=1$ the right side is at
least $1\ge\omega(x,y)$; otherwise $\omega(x',y)=d(x',w)+\rho(w,y)$ with $w\in F$, and
\eqref{eq:glue} gives
\[
\omega(x,y)\;\le\;d(x,w)+\rho(w,y)\;\le\;d(x,x')+d(x',w)+\rho(w,y)\;=\;\omega(x,x')+\omega(x',y).
\]
The inequality $\omega(x',y)\le\omega(x,x')+\omega(x,y)$ is the same statement with $x$ and
$x'$ exchanged. For $\omega(x,x')\le\omega(x,y)+\omega(x',y)$: if either term on the right equals $1$ the
inequality holds since $d\le1$; otherwise $\omega(x,y)=d(x,w)+\rho(w,y)$ and
$\omega(x',y)=d(x',w')+\rho(w',y)$ with $w,w'\in F$, whence
\[
d(x,x')\;\le\;d(x,w)+d(w,w')+d(w',x')\;\le\;d(x,w)+\rho(w,y)+\rho(y,w')+d(w',x'),
\]
again using $d|_F=\rho|_F$ and the triangle inequality for $\rho$ at $[w,y,w']$. The right
side is $\omega(x,y)+\omega(x',y)$.

\smallskip
\noindent\emph{Case 3: $x\in P$, $y,y'\in Q$.}
Symmetric to Case 2.
\end{proof}

\begin{remark}\label{rem:pushout}
The truncation at $1$ in \eqref{eq:glue} is what keeps $\omega$ inside $\bM$; without it
one obtains the usual amalgamated metric, the pushout of $(X,d)\leftarrow(F,d|_F)\to
(Y,\rho)$ in the category of pseudometric spaces and $1$-Lipschitz maps. The dichotomy in
Lemma~\ref{lem:glue} says that each new distance sits at the ceiling or on a geodesic
through $F$, and it is this dichotomy, rather than the amalgamation itself, that the proof
of Theorem~\ref{thm:main} consumes.
\end{remark}

%%%%%%%%%%%%%%%%%%%%%%%%%%%%%%%%%%%%%%%%%%%%%%%%%%%%%%%%%%
\section{Subdividing a distance}
\label{sec:subdivide}
%%%%%%%%%%%%%%%%%%%%%%%%%%%%%%%%%%%%%%%%%%%%%%%%%%%%%%%%%%

No finite extreme metric has an irrational distance, so an irrational distance cannot be
frozen by the method of Section~\ref{sec:proof} directly. It can be written as a sum of
rational distances along a countable geodesic, after which freezing the summands freezes
the sum.

\begin{lemma}\label{lem:subdivide}
Let $(X,d)$ be a bounded-by-$1$ pseudometric space and let $x_0,y_0\in X$ with
$t:=d(x_0,y_0)>0$. There exist a countable set $C=\{c_1,c_2,\ldots\}$ disjoint from $X$ and
a metric $\omega\in\bM(X\cup C)$ with $\omega|_{X\times X}=d$ such that, writing $c_0:=x_0$:
\begin{enumerate}
\item $\omega(c_i,c_{i+1})=a_i$ for a sequence of positive rationals with
      $\sum_{i\ge0}a_i=t$;
\item $\omega(c_0,c_j)=\sum_{i<j}a_i$ and $\omega(c_j,y_0)=t-\sum_{i<j}a_i$ for all
      $j\ge1$;
\item every perturbation $\eps$ of $\omega$ satisfies
      $\displaystyle\eps(x_0,y_0)=\sum_{i\ge0}\eps(c_i,c_{i+1})$, the series converging
      absolutely.
\end{enumerate}
\end{lemma}

\begin{proof}
Choose positive rationals $a_0,a_1,\ldots$ with $\sum_{i\ge0}a_i=t$, put $s_0=0$ and
$s_j=\sum_{i<j}a_i$, so $s_j\uparrow t$ and $0<a_i<t\le1$. Let
$S=\{s_0,s_1,\ldots\}\cup\{t\}\subseteq[0,1]$ carry the metric inherited from $\R$, and let
$C^{+}$ be a copy of $S$ in which the points $s_0$ and $t$ are identified with $x_0$ and
$y_0$ respectively and the points $s_j$ with $j\ge1$ are new; write $c_j$ for the point at
$s_j$ and $C=\{c_1,c_2,\ldots\}$. The metric inherited from $\R$ agrees with $d$ on
$C^{+}\cap X=\{x_0,y_0\}$, since $|t-0|=t=d(x_0,y_0)$, and that intersection is finite, so
Lemma~\ref{lem:glue} supplies $\omega\in\bM(X\cup C)$ restricting to $d$ on $X$ and to the
real-line metric on $C^{+}$. Items (1) and (2) are then immediate.

For (3), let $\eps$ be a perturbation of $\omega$. By (2) the triples $[c_0,c_j,c_{j+1}]$
and $[c_0,c_j,y_0]$ are degenerate for $\omega$, so Lemma~\ref{lem:tight} gives
\[
\eps(c_0,c_{j+1})=\eps(c_0,c_j)+\eps(c_j,c_{j+1}),
\qquad
\eps(c_0,c_j)+\eps(c_j,y_0)=\eps(x_0,y_0),
\]
whence $\eps(c_0,c_j)=\sum_{i<j}\eps(c_i,c_{i+1})$ for every $j$. By
Lemma~\ref{lem:bound}, $|\eps(c_i,c_{i+1})|\le a_i$, so the series converges absolutely,
and $|\eps(c_j,y_0)|\le\omega(c_j,y_0)=t-s_j\to0$. Letting $j\to\infty$ in the second
display gives $\eps(x_0,y_0)=\sum_{i\ge0}\eps(c_i,c_{i+1})$.
\end{proof}

%%%%%%%%%%%%%%%%%%%%%%%%%%%%%%%%%%%%%%%%%%%%%%%%%%%%%%%%%%
\section{Proof of the theorem}
\label{sec:proof}
%%%%%%%%%%%%%%%%%%%%%%%%%%%%%%%%%%%%%%%%%%%%%%%%%%%%%%%%%%

The finite theory enters exactly once.

\begin{proposition}\label{prop:realize}
For every rational $q\in(0,1]$ there are a finite set $F$, an extreme point
$\rho\in\ex(\bM(F))$, and points $u,v\in F$ with $\rho(u,v)=q$.
\end{proposition}

\begin{proof}
Write $q=a/b$ with $b\geq 2$ and $1\leq a\leq b$; this is possible for every rational in
$(0,1]$, doubling numerator and denominator when the reduced denominator is $1$.

Let $B_b$ denote the generalized bowtie metric of \cite{FK-cone} on the cells
$P_1,\ldots,P_{b+1}$ of sizes $1,2,2,\ldots,2$, so that $F$ has $2b+1$ points. Its
interior cells $P_2,\ldots,P_b$ all have size $2$, and there are $b+1\geq 3$ cells, so the
side condition on two-cell bowties is vacuous; and $2b+1\geq 5$. The generalized bowtie
theorem of \cite{FK-cone} therefore makes $B_b$ a generator of an extreme ray of the metric
cone on $F$. Its values on distinct points are the level differences $1,\ldots,b$ together
with the cell-mate value $2$, so $\tfrac{1}{b}B_b$ takes exactly the values
$\tfrac1b,\tfrac2b,\ldots,\tfrac bb$ and has maximum $1$. A generator of an extreme ray with
maximum value $1$ is an extreme point of the body \cite{FK-cone}, and $a/b$ occurs among its
values.
\end{proof}

\begin{remark}\label{rem:bowtie-direct}
The single family $B_b$ covers every rational in $(0,1]$, where the corresponding corollary
of \cite{FK-cone} splits into the cases $q=2$ and $q\geq3$ with different cell patterns.

The formalization of Proposition~\ref{prop:realize} does not pass through the metric cone at
all, and the argument is worth recording, since the final step is shorter in the body than
in the cone. The tight triples of $B_b$ are exactly the \emph{chains}, on points of three
strictly increasing levels, and the \emph{vees}, on two cell-mates and a point of an
adjacent level. Let $\eps$ be a perturbation of $\tfrac1bB_b$. For a four-cycle $x,u,y,v$
whose four sides join consecutive cells and whose two diagonals are realized as vees through
both of the remaining two points, Lemma~\ref{lem:tight} applied to the four vees gives
\[
\eps(x,u)+\eps(u,y)=\eps(x,v)+\eps(v,y),
\qquad
\eps(x,u)+\eps(x,v)=\eps(u,y)+\eps(v,y),
\]
whence $\eps(u,y)=\eps(x,v)$ and $\eps(x,u)=\eps(v,y)$. Propagating these equalities across
and between blocks gives a common value $\lambda$ of $\eps$ on all pairs joining consecutive
cells, and then Lemma~\ref{lem:tight} on vees and chains gives $\eps=\lambda B_b$
throughout. The apex $P_1$ and any point of $P_{b+1}$ are at distance $1$, so
Lemma~\ref{lem:bound} forces $\eps$ to vanish there, and that value is $\lambda b/b=\lambda$.
Hence $\lambda=0$ and $\eps=0$.
\end{remark}

\begin{proof}[Proof of Theorem~\ref{thm:main}]
Fix a countable dense $A\subseteq X$.

\smallskip
\noindent\emph{The construction.}
We build an increasing chain $X=X_0\subseteq X_1\subseteq X_2\subseteq\cdots$ of sets and a
compatible chain of metrics $d=\omega_0,\omega_1,\ldots$ with
$\omega_{k+1}|_{X_k\times X_k}=\omega_k$, each step being one of:

\begin{itemize}
\item[(R)] \emph{Rigidify a rational pair.} Given $u,v\in X_k$ with
$\omega_k(u,v)=q\in\Q\cap(0,1)$, take $F$ and $\rho\in\ex(\bM(F))$ from
Proposition~\ref{prop:realize} realizing $q$ at a pair of points of $F$, identify those two
points of $F$ with $u$ and $v$, and let $X_{k+1}=X_k\cup F$ with $\omega_{k+1}$ the metric
of Lemma~\ref{lem:glue} for the intersection $\{u,v\}$.

\item[(S)] \emph{Subdivide an irrational pair.} Given $u,v\in X_k$ with
$\omega_k(u,v)\notin\Q$, let $X_{k+1}=X_k\cup C$ and $\omega_{k+1}=\omega$ as furnished by
Lemma~\ref{lem:subdivide} applied to $(X_k,\omega_k)$ and the pair $u,v$.
\end{itemize}

Maintain a task list, initially the pairs $\{u,v\}\subseteq A$ with $0<d(u,v)<1$. A task
with rational distance is discharged by a step of type (R); a task with irrational distance
is discharged by a step of type (S), which appends to the list the countably many new pairs
$\{c_i,c_{i+1}\}$, all of rational distance $a_i\in(0,1)$. Each step adds at most countably
many points and at most countably many tasks, so dovetailing arranges that every task is
discharged at some finite stage. Put $\tilde X=\bigcup_k X_k$ and let $\tilde d$ be the
common extension of the $\omega_k$; this is well defined because any two points of
$\tilde X$ lie in a common $X_k$ and the $\omega_k$ are compatible. Then
$\tilde d\in\bM(\tilde X)$, $\tilde d|_{X\times X}=d$, and $\tilde X\setminus X$ is a
countable union of countable sets, hence countable. Since $A\cup(\tilde X\setminus X)$ is
countable and dense in $\tilde X$, the extension is separable.

\smallskip
\noindent\emph{Every perturbation vanishes on $A\times A$.}
Let $\eps$ be a perturbation of $\tilde d$. Let $u,v\in A$ be distinct. If $d(u,v)\in\{0,1\}$
then $\eps(u,v)=0$ by Lemma~\ref{lem:bound}. Otherwise $\{u,v\}$ was a task.

If $d(u,v)$ is rational, the task was discharged by a step of type (R), which placed $u,v$
inside a copy of a finite $F$ carrying an extreme $\rho$. By Corollary~\ref{cor:restrict},
$\eps$ vanishes on $F\times F$, so $\eps(u,v)=0$.

If $d(u,v)$ is irrational, the task was discharged by a step of type (S), and
Lemma~\ref{lem:subdivide}(3) gives $\eps(u,v)=\sum_i\eps(c_i,c_{i+1})$. Each pair
$\{c_i,c_{i+1}\}$ was itself appended to the task list with rational distance $a_i\in(0,1)$
and so was later discharged by a step of type (R); by the previous paragraph
$\eps(c_i,c_{i+1})=0$ for every $i$, and hence $\eps(u,v)=0$.

\smallskip
\noindent\emph{Every perturbation vanishes on $X_k\times X_k$, by induction on $k$.}
For $k=0$: $\eps|_{X\times X}$ is a perturbation of $d$ by Corollary~\ref{cor:restrict}, it
vanishes on $A\times A$ by the previous paragraph, and $A$ is dense in $X$, so
Lemma~\ref{lem:dense} gives $\eps|_{X\times X}=0$.

Assume $\eps$ vanishes on $X_k\times X_k$ and let $X_{k+1}=X_k\cup N$ be the next step,
glued along the finite set $F_0=X_k\cap N$. A pair with both points in $X_k$ is covered by
the inductive hypothesis. For a pair with both points in $N$: if the step was of type (R),
$N$ carries a finite extreme metric and Corollary~\ref{cor:restrict} applies; if the step
was of type (S), $N$ is the chain $\{c_j\}$ together with $F_0=\{u,v\}$, every distance
$\omega_{k+1}(c_i,c_j)$ with $i<j$ is the sum $\sum_{i\le l<j}a_l$ of consecutive ones by
Lemma~\ref{lem:subdivide}(2), so repeated use of Lemma~\ref{lem:tight} expresses
$\eps(c_i,c_j)$ as $\sum_{i\le l<j}\eps(c_l,c_{l+1})$, which vanishes as above; the pairs
$\{c_j,v\}$ are handled by the second display in the proof of
Lemma~\ref{lem:subdivide}(3). Finally let $x\in X_k\setminus N$ and $y\in N\setminus X_k$.
By Lemma~\ref{lem:glue} either $\omega_{k+1}(x,y)=1$, in which case $\eps(x,y)=0$ by
Lemma~\ref{lem:bound}, or there is $z\in F_0$ with
$\omega_{k+1}(x,y)=\omega_{k+1}(x,z)+\omega_{k+1}(z,y)$, in which case
Lemma~\ref{lem:tight} gives $\eps(x,y)=\eps(x,z)+\eps(z,y)$; here $x,z\in X_k$ and
$y,z\in N$, so both summands vanish by what precedes. This completes the induction.

\smallskip
Every pair of points of $\tilde X$ lies in some $X_k$, so $\eps=0$. By
Proposition~\ref{prop:extreme-iff}, $\tilde d$ is extreme.
\end{proof}

\begin{proof}[Proof of Corollary~\ref{cor:irrational}]
Let $V\subseteq[0,1]$ be countable. Realize $V$ as the set of distances of some countable
bounded-by-$1$ pseudometric space---for instance, take a countable set $\{p_v\}_{v\in V}$
together with a base point $p_*$, and define $d(p_*,p_v)=v$ and $d(p_v,p_w)=\max\{v,w\}$
for $v\neq w$, an ultrametric. Apply Theorem~\ref{thm:main}; the extension restricts to
this space, so its distances include $V$. The resulting space is countable, being countable
plus countable.
\end{proof}

%%%%%%%%%%%%%%%%%%%%%%%%%%%%%%%%%%%%%%%%%%%%%%%%%%%%%%%%%%
\section{The canonical extension}
\label{sec:canonical}
%%%%%%%%%%%%%%%%%%%%%%%%%%%%%%%%%%%%%%%%%%%%%%%%%%%%%%%%%%

The construction of Section~\ref{sec:proof} depends on an enumeration of tasks, on the
rationals chosen to subdivide each irrational distance, and on the finite extreme metric
chosen to freeze each rational one. All three choices can be made canonically, and the
enumeration can be eliminated outright, by decorating \emph{ordered pairs} in two
simultaneous stages.

\begin{definition}[canonical selectors]\label{def:selectors}
For irrational $t\in(0,1)$ let $t=\sum_{k\in K(t)}2^{-k}$ be the binary expansion, $K(t)$
infinite, and let the \emph{canonical subdivision} of $t$ be the sequence of gaps
$(2^{-k})_{k\in K(t)}$ in increasing order of $k$. For rational $q\in(0,1)$ let $b(q)$ be
the denominator of $q$ in lowest terms, so $b(q)\geq2$, and let the \emph{canonical piece}
for $q$ be the scaled bowtie $\tfrac{1}{b(q)}B_{b(q)}$ of Proposition~\ref{prop:realize},
with its designated ordered pair (apex, point of level $a$), where $q=a/b(q)$.
\end{definition}

\begin{definition}[canonical extension]\label{def:canonical}
Let $(A,d)$ be a countable metric space. \emph{Stage 1}: for every ordered pair
$(u,v)\in A^2$ with $d(u,v)$ irrational, adjoin the subdivision chain from $u$ to $v$ with
the canonical gaps, gluing the whole family at once. \emph{Stage 2}: for every ordered pair
$(x,y)$ of points of the stage-1 space whose distance is a rational in $(0,1)$ and which
is either a pair from $A$ or an adjacent gap pair of a stage-1 chain, adjoin the canonical
piece for that distance with $x$ at the apex, again gluing the family at once. Write
$\E(A)$ for the result.
\end{definition}

Since each ordered pair is decorated, and $(u,v)$ and $(v,u)$ receive mirror-image
decorations, no orientation is chosen; since the decorations are indexed by the pairs
themselves, no enumeration is chosen; and the selectors of Definition~\ref{def:selectors}
are functions of the distance alone. The proof of Theorem~\ref{thm:main} applies
verbatim---the task list there is exactly the index set here---so $\E(A)$ is an extreme
metric space extending $A$ by countably many points.

\begin{theorem}[functoriality]\label{thm:functor}
$\E$ is a functor from countable bounded-by-$1$ pseudometric spaces with isometric
embeddings to extreme ones: an isometric embedding $\psi\colon A\to B$ extends canonically
to an isometric embedding $\E(\psi)\colon\E(A)\to\E(B)$ restricting to $\psi$, and
$\E(\psi'\circ\psi)=\E(\psi')\circ\E(\psi)$. In particular isometric countable spaces have
isometric canonical extensions.
\end{theorem}

\begin{proof}
$\psi$ carries an ordered pair of $A$ to an ordered pair of $B$ at the same distance, hence
with the same canonical subdivision or the same canonical piece; define $\E(\psi)$ as
$\psi$ on $A$ and as the identity on decoration patterns over matched pairs. Distances
inside $A$ are preserved by hypothesis; distances inside a decoration depend only on the
distance of its pair; and a distance between two different decorations, or between a
decoration and a base point, is by \eqref{eq:glue} a truncated minimum of sums of
within-piece distances and base distances over the finitely many anchors of the pieces
involved, every term of which is matched. The extra decorations present in $\E(B)$ over
unmatched pairs do not enter these formulas. Composition holds because patterns map by the
identity.
\end{proof}

\begin{remark}\label{rem:minimality}
Functoriality and minimality pull apart: a functor must decorate every pair uniformly, so
it cannot fix even the extreme inputs with a distance in $(0,1)$, while a minimal extension
of an extreme space is the space itself. The canonical extension is visibly non-minimal.
Whether minimal extreme extensions exist at all remains open;
see Question~\ref{q:minimal}.
\end{remark}

%%%%%%%%%%%%%%%%%%%%%%%%%%%%%%%%%%%%%%%%%%%%%%%%%%%%%%%%%%
\section{Recovery and nonclassification}
\label{sec:nonclass}
%%%%%%%%%%%%%%%%%%%%%%%%%%%%%%%%%%%%%%%%%%%%%%%%%%%%%%%%%%

Theorem~\ref{thm:main} says that extremality excludes no separable isometry type, and in
this sense already reads as a nonclassification statement. This section makes the reading
precise, in the framework of Borel reducibility \cite{GaoKechris03,Hjorth00}: writing
$\leB$ for Borel reducibility, $\GI$ for graph isomorphism on countable graphs, and
$\isoi$ for isometry of Polish metric spaces, we place isometry of \emph{extreme} Polish
metric spaces above $\GI$.

Extremality passes to completions, so the discussion may be conducted among Polish spaces.

\begin{lemma}\label{lem:completion}
Let $d\in\bM(X)$ be extreme and let $\bar X$ be the completion of $(X,d)$, with metric
$\bar d$. Then $\bar d$ is an extreme point of $\bM(\bar X)$.
\end{lemma}

\begin{proof}
A perturbation of $\bar d$ restricts to a perturbation of $d$ by
Corollary~\ref{cor:restrict}, hence vanishes on $X\times X$; and $X$ is dense in $\bar X$,
so the perturbation vanishes by Lemma~\ref{lem:dense}.
\end{proof}

Extremality is moreover a Borel condition in the standard presentation of Polish metric
spaces as completions of metrics on $\N$: the set $\bM(\N)$ is a compact convex metrizable
subset of $[0,1]^{\N\times\N}$, its extreme points form a $G_\delta$ subset, and by
Lemma~\ref{lem:completion} together with Corollary~\ref{cor:restrict} a metric on $\N$ is
extreme exactly when the metric of its completion is. Write $\isox$ for the restriction of
$\isoi$ to this $G_\delta$ class.

The decorations of the canonical extension are metrically inert, while its chains
accumulate exactly at the core; this makes the core definable from the isometry type.

\begin{lemma}[isolation]\label{lem:isolation}
Every decoration point of $\E(A)$ is an isolated point of $\overline{\E(A)}$.
\end{lemma}

\begin{proof}
Let $p$ be a decoration point, belonging to the piece $P$ glued along the finite anchor set
$F_P$. By \eqref{eq:glue}, the distance from $p$ to any point outside $P$ is at least
$\delta_1:=\min_{z\in F_P}d(p,z)$, a minimum over a finite set of positive numbers. Within
$P$: if $P$ is a bowtie, $P$ is finite and $\delta_2:=$ its least positive distance to $p$
is positive; if $P$ is a chain and $p=c_j$, the within-piece distances from $p$ are
$|s-s_j|$ over the other partial sums $s$ and the two endpoint distances $s_j$ and
$t-s_j$, all at least $\delta_2:=\min\{a_{j-1},a_j,s_j,t-s_j\}>0$. So no point of $\E(A)$
other than $p$ lies within $\min\{\delta_1,\delta_2\}$ of $p$, and the same holds in the
completion.
\end{proof}

\begin{lemma}\label{lem:limits}
Every limit point of $\overline{\E(A)}$ lies in $\bar A$.
\end{lemma}

\begin{proof}
Let $p_n\to z$ with the $p_n$ distinct. If infinitely many $p_n$ lie in $A$ then
$z\in\bar A$ at once, so suppose all are decoration points. If some piece contains
infinitely many $p_n$, then $z$ lies in the closure of that piece, which is the piece
itself when a bowtie, and the piece together with its two endpoints when a chain; a chain
has no limit points but its endpoints, so $z$ is an endpoint, hence a base point of its
stage. If instead the pieces are eventually distinct, then since for $p,p'$ in distinct
pieces \eqref{eq:glue} gives $d(p,p')\ge d(p,F_P)$, the Cauchy property forces
$d(p_n,F_{P_n})\to0$ along a subsequence; choosing anchors $u_n\in F_{P_n}$ with
$d(p_n,u_n)\to0$ gives $u_n\to z$ with the $u_n$ base points of the relevant stage. In
either case $z$ is a limit of stage-2 base points, that is, of points of $A$ together with
chain points; applying the same argument to a sequence of chain points, whose anchors lie
in $A$, puts $z\in\bar A$, which is closed in $\overline{\E(A)}$.
\end{proof}

\begin{theorem}[recovery]\label{thm:recovery}
Suppose every point of $A$ has some point of $A$ at irrational distance. Then the derived
set of $\overline{\E(A)}$ is exactly $\bar A$. Consequently, for countable $A$ and $B$
satisfying the hypothesis,
\[
\overline{\E(A)}\ \cong\ \overline{\E(B)}
\quad\Longrightarrow\quad
\bar A\ \cong\ \bar B .
\]
\end{theorem}

\begin{proof}
By Lemma~\ref{lem:limits} the derived set is contained in $\bar A$. Conversely let
$a\in A$ and pick $a'\in A$ with $d(a,a')$ irrational; the canonical chain for the ordered
pair $(a',a)$ has its partial sums converging to $d(a',a)$, so its points converge to $a$,
and $a$ is a limit point; a general point of $\bar A$ is a limit of points of $A$, and the
derived set is closed. The consequence holds because an isometry of completions carries
derived sets to derived sets.
\end{proof}

The hypothesis is essential only for the containment $\bar A\subseteq{}$derived set;
Lemmas~\ref{lem:isolation} and~\ref{lem:limits} are unconditional.

\begin{theorem}\label{thm:gi}
$\GI\ \leB\ \isox$.
\end{theorem}

\begin{proof}
Fix irrationals $\alpha=1/\sqrt3$ and $\beta=1/\sqrt5$, both in $(\tfrac13,\tfrac23)$. For
a graph $G$ on $\N$ let $X_G$ be $\N$ with $d(m,n)=\alpha$ for edges and $\beta$ for
non-edges, $m\neq n$. Any two values in $(\tfrac13,\tfrac23)$ satisfy all triangle
inequalities, so $X_G$ is a metric space; it is complete and uniformly discrete, so
$\bar X_G=X_G$; every point has points at irrational distance; and since
$\alpha\neq\beta$, a bijection $\N\to\N$ is an isometry $X_G\to X_H$ if and only if it is
an isomorphism $G\to H$.

The reduction is $G\mapsto$ a code of $\overline{\E(X_G)}$: the carrier of $\E(X_G)$ is
$\N$ together with countably many canonically indexed decoration points, enumerated by a
fixed bijection with $\N$, and each distance is given by an explicit formula---a
within-piece value, or a truncated minimum over finitely many anchor routes---so the map is
Borel. The output is extreme by Section~\ref{sec:canonical} and
Lemma~\ref{lem:completion}.

If $G\cong H$ then $X_G\cong X_H$, so $\overline{\E(X_G)}\cong\overline{\E(X_H)}$ by
Theorem~\ref{thm:functor} and Lemma~\ref{lem:completion}. If conversely
$\overline{\E(X_G)}\cong\overline{\E(X_H)}$, then $X_G\cong X_H$ by
Theorem~\ref{thm:recovery}, hence $G\cong H$.
\end{proof}

\begin{corollary}\label{cor:nonclass}
Isometry of extreme Polish metric spaces is not smooth, and is at least as complicated as
the isomorphism relation of any class of countable structures. In particular extreme Polish
metric spaces admit no Borel classification by real-number invariants, nor by countable
structures.
\end{corollary}

\begin{proof}
$\GI$ is universal among isomorphism relations of classes of countable structures and is
not smooth \cite{Hjorth00}; Theorem~\ref{thm:gi} puts $\isox$ above it, and
$\isox\leB\isoi$ by restriction.
\end{proof}

By the theorem of Gao and Kechris \cite{GaoKechris03}, $\isoi$ is Borel bireducible with
the universal orbit equivalence relation induced by Borel actions of Polish groups, so
$\isox$ now sits between $\GI$ and that maximum. It is tempting to close the gap by
feeding $\E$ richer inputs. That route is closed, for a structural reason worth recording.
Call a map $F$ from codes to metric spaces \emph{marked-canonical} if the isometry type of
$F(A)$ depends only on the isometry type of $A$ as a countable metric space---as $\E$ does
by Theorem~\ref{thm:functor}---rather than only on the completion $\bar A$.

\begin{proposition}\label{prop:obstruction}
No marked-canonical $F$ with a recovery property---isometry of outputs implies isometry of
the countable inputs---can Borel-reduce a universal orbit equivalence relation to isometry
of its outputs: any equivalence relation reducible through such a pair of implications
Borel-reduces to $\GI$.
\end{proposition}

\begin{proof}
The two implications together give $z\,\E\,z'\iff A_z\cong A_{z'}$, isometry of countable
metric spaces. A countable metric space is a countable structure in the countable language
$\{R_q:q\in\Q\}$, $R_q(x,y)\iff d(x,y)<q$, with isomorphism equal to isometry, so this
relation Borel-reduces to $\GI$. A universal orbit equivalence relation does not: by the
theory of turbulence \cite{Hjorth00} there are orbit equivalence relations not classifiable
by countable structures, and these reduce to the universal one.
\end{proof}

So the dependence of the extension on its countable dense set is not bookkeeping to be
routed around; it is the exact boundary of the method, and Theorem~\ref{thm:gi} is optimal
for every construction of this shape.

\begin{question}\label{q:universal}
Is $\isox$ Borel bireducible with the universal orbit equivalence relation of Polish group
actions---equivalently, is $\isoi\leB\isox$? By Proposition~\ref{prop:obstruction}, an
affirmative proof must produce extreme extensions whose isometry type depends only on the
completion of the input, or avoid extension constructions altogether. A negative answer
would be equally striking: extremality would then be a genuine simplifying constraint,
placing $\isox$ strictly between the countable-structure maximum and the universal one, a
region where few natural equivalence relations are known to live.
\end{question}

%%%%%%%%%%%%%%%%%%%%%%%%%%%%%%%%%%%%%%%%%%%%%%%%%%%%%%%%%%
\section{Remaining questions}
\label{sec:questions}
%%%%%%%%%%%%%%%%%%%%%%%%%%%%%%%%%%%%%%%%%%%%%%%%%%%%%%%%%%

\begin{question}\label{q:density}
Theorem~\ref{thm:main} spends countability twice: once on the dense set $A$, and once on
the subdivision of each irrational distance. For a bounded-by-$1$ pseudometric space of
density character $\kappa$, does there exist an extreme extension $\tilde X$ with
$|\tilde X\setminus X|\le\kappa$?
\end{question}

The first use of countability appears inessential---Lemma~\ref{lem:dense} holds for a dense
set of any cardinality, and the task list may be indexed by any well-ordering---while the
second is genuine: a single irrational distance requires countably many points to
subdivide, so a space with $\kappa$ irrational distances calls for $\kappa\cdot\aleph_0$
new points, which equals $\kappa$ for infinite $\kappa$. The obstruction to a transfinite
recursion is bookkeeping rather than geometry, but the dovetailing argument above does not
transcribe directly.

\begin{question}\label{q:minimal}
Does every separable bounded-by-$1$ pseudometric space admit a \emph{minimal} extreme
extension, one no proper subspace of which containing $X$ carries an extreme metric
restricting to $d$? Theorem~\ref{thm:functor} answers the functorial half of the canonicity
question; by Remark~\ref{rem:minimality} no functorial extension is minimal, so the two
notions of canonicity diverge and the minimal one remains open.
\end{question}

A related point deserves recording. For finite $X$ the extreme points of $\bM(X)$ are
sparse and rigid, and the difficulty of the subject lies in identifying them. Theorem
\ref{thm:main} shows that for countable $X$ they are instead dense in a strong sense: every
separable bounded-by-$1$ pseudometric space sits inside one. The quantitative content has
migrated entirely into the finite theory, through
Proposition~\ref{prop:realize}---which distances a finite extreme metric can realize---and
that question remains the interesting one.

%%%%%%%%%%%%%%%%%%%%%%%%%%%%%%%%%%%%%%%%%%%%%%%%%%%%%%%%%%
\section{Formal verification}
\label{sec:verify}
%%%%%%%%%%%%%%%%%%%%%%%%%%%%%%%%%%%%%%%%%%%%%%%%%%%%%%%%%%

The results of Sections~\ref{sec:prelim}--\ref{sec:canonical}, and of
Section~\ref{sec:nonclass} apart from the exceptions marked below, have been formalized in
Lean~4 against Mathlib. The development is in \cite{Deposit}; it compiles free of
\texttt{sorry}, and each theorem named below depends only on \texttt{propext},
\texttt{Classical.choice} and \texttt{Quot.sound}. The formalization of
Section~\ref{sec:nonclass} constructs no completions: isolation, limits and recovery are
proved on the countable carriers, with isolation carrying an explicit radius, and
Theorem~\ref{thm:gi}'s equivalence is proved there for simple graphs, which is the
generality its statement intends. Two entries are marked accordingly:
$^{\dagger}$Lemma~\ref{lem:completion} is the one numbered result left in prose---its
two-line proof consumes only Corollary~\ref{cor:restrict} and Lemma~\ref{lem:dense}, both
formalized---and $^{\ddagger}$the containment of limit points in $\bar A$ is formalized in
the direction the reduction uses. The descriptive-set-theoretic inputs rest on cited
literature outside the scope of formalization, as marked.

Extremality is formalized twice: as the vanishing of every perturbation, and as membership
in Mathlib's \texttt{Set.extremePoints} for the convex set of bounded-by-$1$ pseudometrics,
the two being identified by Proposition~\ref{prop:extreme-iff}. Theorem~\ref{thm:main} and
Corollary~\ref{cor:irrational} carry no hypotheses beyond those stated here.

\begin{center}
\begin{tabular}{lll}
\textbf{Result} & \textbf{Status} & \textbf{Lean name}\\
\hline
Theorem~\ref{thm:main} & formalized & \texttt{separable\_extends\_to\_extreme'}\\
 & & \texttt{separable\_extends\_to\_extremePoint'}\\
Corollary~\ref{cor:irrational} & formalized & \texttt{exists\_countable\_extreme\_realizing'}\\
Definition~\ref{def:perturbation} & formalized & \texttt{IsPerturbation}\\
Proposition~\ref{prop:extreme-iff} & formalized & \texttt{rigid\_iff\_extremePoint}\\
Lemma~\ref{lem:bound} & formalized & \texttt{abs\_pert\_le}\\
Lemma~\ref{lem:tight} & formalized & \texttt{pert\_add\_of\_tight}\\
Lemma~\ref{lem:dense} & formalized & \texttt{pert\_eq\_zero\_of\_dense}\\
Corollary~\ref{cor:restrict} & formalized & \texttt{pert\_eq\_zero\_on\_of\_rigid\_restrict}\\
Lemma~\ref{lem:glue} & formalized & \texttt{glue\_isBddPseudo}, \texttt{glue\_dichotomy}\\
Remark~\ref{rem:pushout} & formalized & \texttt{glue\_no\_truncation\_not\_isBddPseudo}\\
Lemma~\ref{lem:subdivide} & formalized & \texttt{pert\_eq\_tsum\_of\_geodesic}\\
Proposition~\ref{prop:realize} & formalized & \texttt{exists\_finite\_rigid\_realizing}\\
Remark~\ref{rem:bowtie-direct} & formalized & \texttt{Bowtie.bowtie\_rigid}\\
\S\ref{sec:proof} induction & formalized & \texttt{rigid\_of\_chain}\\
Theorem~\ref{thm:functor} & formalized & \texttt{E\_functor}, \texttt{E\_functor\_surj}\\
Lemma~\ref{lem:completion} & prose only$^{\dagger}$ & ---\\
Lemmas~\ref{lem:isolation}, \ref{lem:limits} & formalized$^{\ddagger}$ & \texttt{decor\_isolated}, \texttt{core\_is\_limit}\\
Theorem~\ref{thm:recovery} & formalized, carrier form & \texttt{recovery}\\
Theorem~\ref{thm:gi}, metric content & formalized & \texttt{main\_equivalence}\\
Theorem~\ref{thm:gi}, Borel wrapper & cites \cite{GaoKechris03,Hjorth00} & ---\\
Corollary~\ref{cor:nonclass}, Prop.~\ref{prop:obstruction} & cite \cite{GaoKechris03,Hjorth00} & ---\\
Questions~\ref{q:density}, \ref{q:minimal}, \ref{q:universal} & open questions & ---\\
\end{tabular}
\end{center}

Proposition~\ref{prop:realize} is formalized by the direct argument of
Remark~\ref{rem:bowtie-direct} rather than by the route through the metric cone given in its
proof above, so the machine-checked development depends on no result outside this paper.

Four statements were recorded in the development as deliberately false, to guard hypotheses
that the proofs consume: the amalgam of Lemma~\ref{lem:glue} without its truncation at $1$;
the dichotomy of Lemma~\ref{lem:glue} with an infinite gluing locus; the sufficiency of a
single round of gluing, without the subdivision of Lemma~\ref{lem:subdivide}; and the
conclusion of Lemma~\ref{lem:dense} without the hypothesis that $\eps$ is a perturbation.
Each is refuted in the development by an explicit counterexample.

\subsection*{Acknowledgment}
We thank Dmitri Nikshych for the observation that Theorem~\ref{thm:main} should be read as
a nonclassification result; Sections~\ref{sec:canonical} and~\ref{sec:nonclass} grew from
pursuing it.

\begin{thebibliography}{9}

\bibitem{Avis80}
D.~Avis, On the extreme rays of the metric cone,
\textit{Canad.\ J.\ Math.}~\textbf{32} (1980), no.~1, 126--144.

\bibitem{Bandelt92}
H.-J.~Bandelt and A.W.M.~Dress, A canonical decomposition theory for metrics on a finite
set, \textit{Adv.\ Math.}~\textbf{92} (1992), no.~1, 47--105.

\bibitem{Deposit}
D.V.~Feldman and E.R.~Kehoe, \textit{Every separable metric space extends to an extreme
one: paper and Lean~4 formalization}, Zenodo (2026),
\url{https://doi.org/10.5281/zenodo.21929312}.

\bibitem{Deza97}
M.M.~Deza and M.~Laurent, \textit{Geometry of Cuts and Metrics}, Springer-Verlag,
Berlin--Heidelberg, 1997.

\bibitem{GaoKechris03}
S.~Gao and A.S.~Kechris, \textit{On the classification of Polish metric spaces up to
isometry}, Mem.\ Amer.\ Math.\ Soc.\ \textbf{161} (2003), no.~766.

\bibitem{Hjorth00}
G.~Hjorth, \textit{Classification and Orbit Equivalence Relations}, Math.\ Surveys
Monogr.~75, Amer.\ Math.\ Soc., 2000.

\bibitem{FK-cone}
D.V.~Feldman and E.R.~Kehoe, Extreme rays of the metric cone and the metric body: bowtie
metrics and half-one extremality, companion paper; the results quoted here also appear as
\cite[Ch.~2]{Kehoe19}.

\bibitem{Kehoe19}
E.R.~Kehoe, \textit{Pseudometrics, the Complex of Ultrametrics, and Iterated Cycle
Structures}, Ph.D.\ dissertation, University of New Hampshire, 2019.
\url{https://scholars.unh.edu/dissertation/2451}

\bibitem{Koolen00}
J.~Koolen, V.~Moulton, and U.~T\"{o}nges, A classification of the six-point prime metrics,
\textit{European J.\ Combin.}~\textbf{21} (2000), no.~6, 815--829.

\end{thebibliography}

\end{document}
